\section{Conclusiones}
Una vez finalizado el desarrollo de la aplicación, resulta productivo dedicar un espacio a la reflexión del trabajo realizado y para evaluar si los objetivos planteados han sido cumplidos satisfactoriamente.

En lo que concierne a los objetivos funcionales, todos ellos han podido ser implementados satisfactoriamente, aunque no todos han llevado el esfuerzo esperado originalmente. Por ejemplo, la visualización de las actividades en un mapa ha sido sorprendentemente sencilla, mientras que otras funcionalidades, como gestionar la visibilidad de actividades compartidas en murales, han resultado ser sorprendentemente complejas. También ha habido funcionalidades que han resultado tener limitaciones distintas a las esperadas. Por ejemplo, se había asumido que un prerrequisito para poder comparar dos actividades sería que hayan ocurrido en la misma ruta, pero ha sido posible realizar comparaciones entre actividades en distintas rutas sin dificultad añadida, ampliando el alcance del objetivo sin coste adicional. 

Por otro lado, en lo referente a los objetivos técnicos, estos también han sido cumplidos. Y en el proceso, se ha obtenido mucha experiencia con las distintas tecnologías y se han conocido herramientas adicionales ofrecidas por los frameworks y tecnologías. En particular, al realizar un trabajo más profesional con Spring y Angular se han descubierto varias herramientas y procesos que, al no haber sido necesarios en proyectos realizados durante la carrera, le eran desconocidos al autor. Ejemplos de esto son el uso de un interceptor HTTP en el cliente, o un gestor de excepciones global en el servidor. En cuanto a complejidad, el despliegue de la aplicación en Azure ha sido el objetivo técnico más difícil de cumplir, al ser una tecnología vasta y compleja con la que el autor no tiene mucha experiencia.

En un nivel más personal, el trabajo realizado durante este TFG ha resultado ser una excelente experiencia, al ofrecer una oportunidad de gestionar un proyecto de software grande y trabajar con tecnologías utilizadas en producción en proyectos reales, en particular Spring y Angular, que han resultado ser las tecnologías empleadas en el proyecto al que estoy el autor ha sido asignado en el trabajo. El proceso de investigación de las distintas tecnologías, el encontrar y solucionar problemas clásicos propios de cada tecnología y framework, y el propio trabajo día a día con ellas aportan una experiencia esencial imposible de obtener de otra manera.

Más allá de eso, el autor se considera satisfecho con la aplicación desarrollada. Si bien tiene mucho espacio para mejorar y para añadir nuevas funcionalidades, considera que cumple los requisitos que tenía para ella al plantearse realizarla.

\section{Trabajos futuros} 
Todo proyecto puede ser mejorado o ampliado, y este no es excepción. Hay infinidad de formas de ampliar el proyecto, y se han planteado inicialmente muchas funcionalidades que al final no han sido desarrolladas, así como funcionalidades planteadas durante el desarrollo, que, para mantener la planificación, han sido dejadas para un futuro desarrollo. Algunas de estas funcionalidades serían: 
\begin{itemize}
    \item Mejorar la pantalla de ruta, añadiendo datos cumulativos, mejor actividad en la ruta, visualización de récords, comparación de una actividad con la media de la ruta etc.
    \item Enviar una invitación para unirse a un mural a través de la aplicación.
    \item Ampliar los elementos mostrados en el dashboard de usuario y Mural, y ofrecer funcionalidades de personalización de los mismos, pudiendo cambiar qué elementos mostrar y en qué orden.
    \item Acceder a las funcionalidades de análisis y comparación de actividades, así como creación y gestión de rutas, sin iniciar sesión ni crear cuenta.
    \item Mostrar en la pantalla de una actividad los récords superados.
    \item Crear récords personalizados.
\end{itemize}

Y como estas, muchas otras funcionalidades que se ha decidido no implementar en esta versión de la aplicación.


En definitiva, aunque la aplicación cumple los objetivos originales, puede ser ampliada en muchas direcciones distintas para formar una aplicación más madura que tenga una oportunidad real de competir con otras aplicaciones en su ámbito.
